% TOP_PROBLEM: {"id": 3000039, "title": "Goddyn's conjecture on thin spanning trees", "category_id": 2, "category": {"id": 2, "name": "combinatorics", "display_name": "Combinatorics", "description": "Counting problems, graph theory, discrete structures.", "slug": "combinatorics", "order_index": 2, "created_at": "2026-07-31T15:26:25.671Z"}, "status": "partially_solved", "problem_number": "AMR-029-0039", "set_id": 13}
% FIRST_POSTED: 2026-10-01
% ENTRY_KIND: statement_only
% UnsolvedMath ID 3000039; source catalogue code AMR-029-0039
% !TeX program = lualatex
% Definition and sources only; this file is not an LLM solution attempt.
\documentclass[11pt,a4paper]{article}
\usepackage[margin=25mm]{geometry}
\usepackage{fontspec}
\setmainfont{lmroman10-regular.otf}[BoldFont=lmroman10-bold.otf,
 ItalicFont=lmroman10-italic.otf,BoldItalicFont=lmroman10-bolditalic.otf]
\usepackage{amsmath,amssymb,mathtools}
\usepackage{unicode-math}
\setmathfont{latinmodern-math.otf}
\AtBeginDocument{\let\setminus\smallsetminus}
\usepackage{xurl,hyperref}
\hypersetup{colorlinks=true,urlcolor=blue}
\setcounter{secnumdepth}{0}
\setlength{\parindent}{0pt}
\setlength{\parskip}{5pt}
\setlength{\emergencystretch}{3em}
\begin{document}
\section*{Goddyn's conjecture on thin spanning trees}\label{3000039}
{\small UnsolvedMath ID 3000039 · AMR-029-0039 · Combinatorics · AMR Open Problem Lists}
\subsection{Definitions and mathematical statement}
Let $G=(V,E)$ be a finite loopless undirected multigraph. For
$X\subseteq V$, let $\delta_G(X)$ be the set of edges with exactly
one endpoint in $X$, counting parallel copies separately. The graph
is $k$-edge-connected if $|\delta_G(X)|\ge k$ for every nonempty proper
subset $X$.

A spanning tree $T\subseteq G$ is $\varepsilon$-thin when
\[
                         |\delta_T(X)|\le\varepsilon|\delta_G(X)|
                                            \quad(X\subseteq V).
\]
Goddyn's conjecture asks whether there exists a function
$f:(0,1)\to\mathbb Z_{>0}$ such that, for each $\varepsilon\in(0,1)$,
every $f(\varepsilon)$-edge-connected graph contains an
$\varepsilon$-thin spanning tree.

The connectivity threshold depends only on the desired thinness and
must be independent of the number of vertices, maximum degree,
embedding genus, and all other graph parameters. The same tree must
satisfy the inequality for all cuts simultaneously; choosing a different
tree for each cut is insufficient. This is a cut-count condition and
does not assume or demand a stronger spectral matrix inequality.
The conjecture requests existence of some threshold function, without
prescribing its rate of growth. A particular construction with a
threshold that also depends on $|V|$ does not establish the stated
uniform conclusion.
\subsection{Short English statement}
Does sufficiently high edge connectivity force a spanning tree
using only a prescribed small fraction of the edges of every cut?
\subsection{Sources}
\begin{itemize}
\item[S1] ULAM.AI and the UnsolvedMath Contributors, supplied problems.json, record 3000039 (AMR-029-0039), AMR Open Problem Lists. Catalogue identity and source transcription; CC BY 4.0.
\url{https://huggingface.co/datasets/ulamai/UnsolvedMath}.
\item[S2] Egres Open - List of open problems, Open-problem page: Goddyn's conjecture on thin spanning trees. Original source indexed by AMR-029-0039.
\url{https://oldlemon.cs.elte.hu/egres/open/List_of_open_problems}.
\item[S3] EGRES Open, Goddyn's conjecture on thin spanning trees. Individual source page and explanatory remarks.
\url{https://lemon.cs.elte.hu/egres/open/Goddyn%27s_conjecture_on_thin_spanning_trees}.
\end{itemize}
\end{document}
